% 图 4-6 第二稿：遵循 nature-figure 矢量示意图规范
%  - 一种字族、两种字号：层名 \small 加粗（唯一加粗处），说明与工具 \footnotesize 常规
%  - 全部文字同一深色；唯一强调色只标“工具类别”这一本图要传达的信息
%  - 同一角色的框等宽、上下边对齐；每个框内部顺序相同：层名 → 一行说明 → （底部）工具类别
%  - 连线标签与其他文字保持间距；虚线只表示一种含义：离线
\begin{figure}[htbp]
\centering
\begin{tikzpicture}[fig, x=1mm, y=1mm]
  % \layerbox{x0}{x1}{ytop}{h}{填充}{边框样式}{层名}{说明}{工具类别，可空}
  \newcommand{\layerbox}[9]{%
    \begin{scope}
      \path[clip, rounded corners=1.5pt] (#1,#3) rectangle (#2,#3-#4);
      \fill[#5] (#1,#3) rectangle (#2,#3-#4);
      \if\relax\detokenize{#9}\relax\else
        \fill[accentpale] (#1,#3-#4) rectangle (#2,#3-#4+5);
        \draw[accent, line width=0.4pt] (#1,#3-#4+5) -- (#2,#3-#4+5);
      \fi
    \end{scope}
    \draw[figline, rounded corners=1.5pt, #6] (#1,#3) rectangle (#2,#3-#4);
    \if\relax\detokenize{#9}\relax
      \node[align=center, text=ink] at ({(#1+#2)/2},{#3-#4/2})
        {{\small\bfseries #7}\\[0.5pt]{\footnotesize #8}};
    \else
      \node[align=center, text=ink] at ({(#1+#2)/2},{#3-(#4-5)/2})
        {{\small\bfseries #7}\\[0.5pt]{\footnotesize #8}};
      \node[text=ink, font=\figfont\footnotesize] at ({(#1+#2)/2},{#3-#4+2.5}) {工具：#9};
    \fi
  }
  % ---- 第 1 层 ----
  \layerbox{-55}{55}{0}{11}{figfill}{}{接口层}{网页、移动端、开放 API}{}
  % ---- 第 2 层 ----
  \layerbox{-55}{55}{-19}{16}{figfill}{}{编排层}{工作流与智能体、提示词组织、权限校验、安全审查、日志}{编排平台、应用开发框架}
  % ---- 第 3 层：三个等宽框，等间距 ----
  \layerbox{-55}{-23}{-48}{16}{figfill}{}{知识层}{向量检索、搜索、文档库}{向量数据库}
  \layerbox{-16}{16}{-48}{16}{figfill}{}{模型服务层}{推理引擎、模型 API}{部署工具}
  \layerbox{23}{55}{-48}{16}{white}{dashed}{微调}{离线训练，产出新权重}{微调工具}
  % ---- 第 4 层 ----
  \layerbox{-55}{55}{-76}{11}{white}{}{基础设施层}{计算芯片、存储、互联、驱动与软件栈（第\,4.2\,节）}{}
  % ---- 连线：主路径实线，离线虚线 ----
  \draw[<->, line width=0.6pt] (0,-11) -- (0,-19);
  \draw[<->, line width=0.6pt] (-39,-35) -- (-39,-48) node[midway, right=1.5mm, note] {检索资料};
  \draw[<->, line width=0.6pt] (0,-35)   -- (0,-48)   node[midway, right=1.5mm, note] {生成回答};
  \draw[->, dashed, line width=0.6pt] (23,-56) -- (16,-56);
  \foreach \x in {-39,0,39} \draw[->, line width=0.6pt] (\x,-64) -- (\x,-76);
\end{tikzpicture}
\caption{大模型应用的典型分层与五类工具的位置。每个框底部的色带标出该层对应的工具类别。编排层发起检索与模型调用；微调（虚线框）在线下进行，产出的权重交给模型服务层加载}
\label{fig:app-arch}
\end{figure}
